\chapter{§3.7 勾股定理}
\label{sec:3.7}


\prerequisite{§3.4, §2.2}

\gradelevel{8 年级}


\section*{勾股定理的发现}


\begin{explore}


在方格纸上画一个直角三角形，两条直角边分别占 $3$ 格和 $4$ 格。然后在每条边上向外画一个正方形：

\begin{itemize}
  \item 短直角边的正方形面积 $= 3^2 = 9$
  \item 长直角边的正方形面积 $= 4^2 = 16$
  \item 斜边的正方形面积 $= ?$
\end{itemize}


用数格子的方法（或割补法）数一数，你会发现斜边上的正方形面积恰好是 $25$ 。

\[
9 + 16 = 25
\]


也就是 $3^2 + 4^2 = 5^2$ 。巧合吗？换一个直角三角形试试—— $5, 12, 13$ ：

\[
5^2 + 12^2 = 25 + 144 = 169 = 13^2
\]


又成立了！这个规律适用于所有直角三角形吗？

\end{explore}

\begin{understand}


\textbf{勾股定理}（Pythagorean Theorem）：

直角三角形两条直角边的平方和等于斜边的平方。

如果直角三角形的两条直角边为 $a$ 、 $b$ ，斜边为 $c$ ，则：

\[
a^2 + b^2 = c^2
\]


\begin{figure}[htbp]
\centering
\begin{tikzpicture}[every node/.style={font=\small}, scale=0.85]
  % Right triangle: C at origin, A at (4,0), B at (0,3)
  \coordinate (C) at (0,0);
  \coordinate (A) at (4,0);
  \coordinate (B) at (0,3);
  \draw[main line] (A)--(B)--(C)--cycle;
  \node[below right] at (A) {$A$};
  \node[above left]  at (B) {$B$};
  \node[below left]  at (C) {$C$};
  % Right angle mark at C
  \draw (0.3,0)--(0.3,0.3)--(0,0.3);
  % Side labels
  \node[below] at (2,0) {$a = 4$};
  \node[left]  at (0,1.5) {$b = 3$};
  \node[above right, rotate={atan2(3,4)}] at (2,1.5) {$c = 5$};
  % Square on leg a (below)
  \draw[aux line] (0,-4)--(4,-4)--(4,0);
  \draw[aux line] (0,0)--(0,-4);
  \node at (2,-2) {$a^2 = 16$};
  % Square on leg b (left)
  \draw[aux line] (-3,0)--(-3,3)--(0,3);
  \draw[aux line] (0,0)--(-3,0);
  \node at (-1.5,1.5) {$b^2 = 9$};
  % Label for hypotenuse square
  \node[right=0.3cm] at (4,1.5) {$c^2 = a^2 + b^2$};
  \node[right=0.3cm] at (4,0.8) {$= 16 + 9 = 25$};
\end{tikzpicture}
\caption{勾股定理：直角三角形中 $a^2 + b^2 = c^2$（本例：$3^2+4^2=5^2$）}
\label{fig:pythagorean-theorem}
\end{figure}


这个定理在中国古代被称为"勾股定理"——直角三角形的短直角边叫"勾"，长直角边叫"股"，斜边叫"弦"。《周髀算经》中记载："勾三股四弦五"，说的就是 $3^2 + 4^2 = 5^2$ 。

\end{understand}

\begin{quote}
在西方，这个定理以古希腊数学家毕达哥拉斯（Pythagoras）的名字命名。但中国、巴比伦、印度等古文明都独立发现了这一定理。
\end{quote}


\begin{keytakeaway}


\begin{itemize}
  \item 勾股定理： $a^2 + b^2 = c^2$ （ $c$ 为斜边）。
  \item 只适用于\textbf{直角三角形}。
  \item 已知任意两边，可以求第三边。
\end{itemize}


\end{keytakeaway}



\section*{勾股定理的证明}


\begin{explore}


勾股定理有数百种证明方法。这里介绍一种最直观的"拼正方形"证明法。

\end{explore}

\begin{understand}


\textbf{证明（赵爽弦图法）}：

用四个完全相同的直角三角形（直角边为 $a$ 、 $b$ ，斜边为 $c$ ）拼成一个大正方形。

\begin{figure}[htbp]
\centering
\begin{tikzpicture}[every node/.style={font=\small}, scale=0.85]
  % Outer square, side c = 5
  \def\c{5} \def\a{3} \def\b{4}
  % Outer square corners
  \draw[main line] (0,0)--(\c,0)--(\c,\c)--(0,\c)--cycle;
  % Four inner points dividing each side into a and b
  \coordinate (P1) at (\a,0);
  \coordinate (P2) at (\c,\a);
  \coordinate (P3) at (\b,\c);
  \coordinate (P4) at (0,\b);
  % Inner tilted square (c×c tilted)
  \draw[main line, draw=blue!70!black] (P1)--(P2)--(P3)--(P4)--cycle;
  \fill[blue!10] (P1)--(P2)--(P3)--(P4)--cycle;
  % Labels on outer square sides
  \node[below] at (\a/2,0) {$a$};
  \node[below] at ({(\a+\c)/2},0) {$b$};
  \node[blue!70!black] at (\c/2+0.1,\c/2) {$c^2$};
  % Labels on two triangles
  \node[font=\scriptsize] at (0.6,0.8) {$\frac{1}{2}ab$};
  \node[font=\scriptsize] at (\c-0.5,\c/2) {$\frac{1}{2}ab$};
\end{tikzpicture}
\caption{赵爽弦图证明：大正方形面积 $= c^2 + 4 \times \frac{1}{2}ab$，同时 $= (a+b)^2$，整理得 $c^2 = a^2+b^2$}
\label{fig:zhao-shuang-proof}
\end{figure}


将四个直角三角形排列如下：大正方形的外边长为 $c$ ，内部留出一个小正方形，边长为 $\lvert b - a \rvert$ 。

大正方形的面积 $= c^2$ 。

大正方形面积也等于四个三角形面积 $+$ 内部小正方形面积：

\[
c^2 = 4 \times \frac{1}{2}ab + (b - a)^2
\]


\[
c^2 = 2ab + b^2 - 2ab + a^2
\]


\[
c^2 = a^2 + b^2
\]


证毕。

\textbf{另一种理解}：也可以将四个相同直角三角形拼成一个边长为 $(a + b)$ 的大正方形。

$(a + b)^2 = c^2 + 4 \times \dfrac{1}{2}ab$

$a^2 + 2ab + b^2 = c^2 + 2ab$

$a^2 + b^2 = c^2$

\end{understand}

\begin{keytakeaway}


\begin{itemize}
  \item 赵爽弦图法是最经典的中国式证明。
  \item 核心思想是面积相等：用两种方式计算同一图形的面积。
\end{itemize}


\end{keytakeaway}



\section*{勾股定理的应用}


\begin{explore}


一架 $5$ 米长的梯子斜靠在墙上，梯脚距墙 $3$ 米。梯子顶端离地面多高？

\end{explore}

\begin{understand}


勾股定理的直接应用：\textbf{已知直角三角形两边，求第三边。}

\begin{itemize}
  \item 已知两直角边 $a$ 、 $b$ ，求斜边： $c = \sqrt{a^2 + b^2}$
  \item 已知斜边 $c$ 和一条直角边 $a$ ，求另一条直角边： $b = \sqrt{c^2 - a^2}$
\end{itemize}


\end{understand}

\begin{workedexamples}


\textbf{例 1.} 一架 $5$ 米长的梯子斜靠在墙上，梯脚距墙 $3$ 米。求梯子顶端离地面的高度。

\textbf{解：}
墙、地面和梯子构成直角三角形，梯子是斜边。

设梯顶离地高度为 $h$ ，则 $h^2 + 3^2 = 5^2$ 。

$h^2 = 25 - 9 = 16$ ， $h = 4$ 米。

\textbf{例 2.} 直角三角形的两条直角边分别为 $6$ 和 $8$ ，求斜边和面积。

\textbf{解：}
斜边 $c = \sqrt{6^2 + 8^2} = \sqrt{36 + 64} = \sqrt{100} = 10$ 。

面积 $= \dfrac{1}{2} \times 6 \times 8 = 24$ 。

\textbf{例 3.} 一块长方形草地长 $40$ m，宽 $30$ m。小明从一个角走到对角，走了多少米？

\textbf{解：}
对角线就是斜边： $d = \sqrt{40^2 + 30^2} = \sqrt{1600 + 900} = \sqrt{2500} = 50$ m。

\textbf{例 4.} 在等腰三角形 $ABC$ 中， $AB = AC = 10$ cm， $BC = 12$ cm。求 $BC$ 边上的高 $AD$ 的长度。

\textbf{解：}
因为等腰三角形底边上的高平分底边， $BD = \dfrac{1}{2}BC = 6$ cm。

在直角 $\triangle ABD$ 中， $AD = \sqrt{AB^2 - BD^2} = \sqrt{100 - 36} = \sqrt{64} = 8$ cm。

\end{workedexamples}

\begin{keytakeaway}


\begin{itemize}
  \item 勾股定理的核心应用：已知两边求第三边。
  \item 实际问题中，寻找或构造直角三角形是关键步骤。
\end{itemize}


\end{keytakeaway}

\begin{practice}


\begin{enumerate}
  \item 直角三角形的两条直角边分别为 $5$ cm 和 $12$ cm，求斜边。
  \item 直角三角形的斜边为 $17$ cm，一条直角边为 $8$ cm，求另一条直角边。
  \item 一棵树在距地面 $3$ m 处折断，树尖落在地面上距树根 $4$ m 处。求树折断前的高度。
  \item 正方形的对角线长为 $8$ cm，求边长。
\end{enumerate}


\end{practice}



\section*{勾股定理的逆定理}


\begin{explore}


勾股定理说"如果是直角三角形，那么 $a^2 + b^2 = c^2$ "。反过来能成立吗？如果三条边满足 $a^2 + b^2 = c^2$ ，能否断定它是直角三角形？

取三条线段，长分别为 $5, 12, 13$ 。检验： $5^2 + 12^2 = 25 + 144 = 169 = 13^2$ 。用这三条边围成三角形——用量角器量一量最大角，确实是 $90°$ ！

\end{explore}

\begin{understand}


\textbf{勾股定理的逆定理}：如果三角形的三边长 $a$ 、 $b$ 、 $c$ 满足 $a^2 + b^2 = c^2$ ，则这个三角形是直角三角形（ $c$ 所对的角为直角）。

\end{understand}

\begin{quote}
能使 $a^2 + b^2 = c^2$ 成立的三个正整数 $(a, b, c)$ 叫做\textbf{勾股数}。常见的勾股数组有：
- $(3, 4, 5)$ 及其倍数 $(6, 8, 10)$ 、 $(9, 12, 15)$ 、...
- $(5, 12, 13)$
- $(8, 15, 17)$
- $(7, 24, 25)$
\end{quote}


\textbf{推广应用——判断三角形的类型}：

设三角形三边 $a \leq b \leq c$ ：
\begin{itemize}
  \item $a^2 + b^2 = c^2$ ：直角三角形
  \item $a^2 + b^2 > c^2$ ：锐角三角形
  \item $a^2 + b^2 < c^2$ ：钝角三角形
\end{itemize}


\begin{workedexamples}


\textbf{例 1.} 判断下列各组数能否构成直角三角形的三边：
(a) $6, 8, 10$ ；(b) $7, 8, 11$ ；(c) $9, 40, 41$ 。

\textbf{解：}
(a) $6^2 + 8^2 = 36 + 64 = 100 = 10^2$ 。是直角三角形。

(b) $7^2 + 8^2 = 49 + 64 = 113 \neq 121 = 11^2$ 。且 $113 < 121$ ，所以 $7^2 + 8^2 < 11^2$ ，是钝角三角形。

(c) $9^2 + 40^2 = 81 + 1600 = 1681 = 41^2$ 。是直角三角形。

\textbf{例 2.} $\triangle ABC$ 中， $AB = 15$ ， $BC = 20$ ， $AC = 25$ 。判断 $\triangle ABC$ 的形状，并求面积。

\textbf{解：}
$AB^2 + BC^2 = 225 + 400 = 625 = 25^2 = AC^2$ 。

所以 $\triangle ABC$ 是直角三角形， $\angle B = 90°$ 。

面积 $= \dfrac{1}{2} \times AB \times BC = \dfrac{1}{2} \times 15 \times 20 = 150$ 。

\textbf{例 3.} 三角形的三边长为 $a = \sqrt{2}$ 、 $b = \sqrt{3}$ 、 $c = \sqrt{5}$ 。判断是否为直角三角形。

\textbf{解：}
$a^2 + b^2 = 2 + 3 = 5 = c^2$ 。

所以这是直角三角形。

\end{workedexamples}

\begin{keytakeaway}


\begin{itemize}
  \item 勾股定理的逆定理可用于判定直角三角形。
  \item 利用 $a^2 + b^2$ 与 $c^2$ 的比较，还能判断三角形是锐角还是钝角。
  \item 常见勾股数： $(3,4,5)$ 、 $(5,12,13)$ 、 $(8,15,17)$ 、 $(7,24,25)$ 。
\end{itemize}


\end{keytakeaway}

\begin{practice}


\begin{enumerate}
  \item 判断以下各组长度能否构成直角三角形：(a) $11, 60, 61$ ；(b) $1, 1, \sqrt{2}$ ；(c) $4, 5, 6$ 。
  \item 一块三角形地，三边长分别为 $9$ m、 $12$ m、 $15$ m。这块地是什么形状？求面积。
  \item 判断三边为 $\sqrt{3}, 2, \sqrt{7}$ 的三角形是什么类型的三角形。
  \item 在 $\triangle ABC$ 中， $a = 2k$ ， $b = k^2 - 1$ ， $c = k^2 + 1$ （ $k > 1$ ）。证明 $\triangle ABC$ 是直角三角形。
\end{enumerate}


\end{practice}



\begin{answer}


\begin{enumerate}
  \item $c = \sqrt{5^2 + 12^2} = \sqrt{25 + 144} = \sqrt{169} = 13$ cm。
\end{enumerate}


\begin{enumerate}
  \item $b = \sqrt{17^2 - 8^2} = \sqrt{289 - 64} = \sqrt{225} = 15$ cm。
\end{enumerate}


\begin{enumerate}
  \item 折断处到地面 $3$ m，折断的部分为斜边。折断部分长 $= \sqrt{3^2 + 4^2} = 5$ m。树折断前的高度 $= 3 + 5 = 8$ m。
\end{enumerate}


\begin{enumerate}
  \item 设边长为 $a$ ，对角线 $= a\sqrt{2} = 8$ ， $a = \dfrac{8}{\sqrt{2}} = 4\sqrt{2} \approx 5.66$ cm。
\end{enumerate}


5.
   (a) $11^2 + 60^2 = 121 + 3600 = 3721 = 61^2$ ，是直角三角形。
   (b) $1^2 + 1^2 = 2 = (\sqrt{2})^2$ ，是直角三角形。
   (c) $4^2 + 5^2 = 16 + 25 = 41 \neq 36 = 6^2$ 。因为 $41 > 36$ ，即 $a^2 + b^2 > c^2$ ，是锐角三角形。

\begin{enumerate}
  \item $9^2 + 12^2 = 81 + 144 = 225 = 15^2$ ，是直角三角形。面积 $= \dfrac{1}{2} \times 9 \times 12 = 54$ m $^2$ 。
\end{enumerate}


\begin{enumerate}
  \item $(\sqrt{3})^2 + 2^2 = 3 + 4 = 7 = (\sqrt{7})^2$ ，是直角三角形。
\end{enumerate}


\begin{enumerate}
  \item 验证 $a^2 + b^2 = c^2$ ：
\end{enumerate}

   $a^2 + b^2 = (2k)^2 + (k^2 - 1)^2 = 4k^2 + k^4 - 2k^2 + 1 = k^4 + 2k^2 + 1 = (k^2 + 1)^2 = c^2$ 。
   所以 $a^2 + b^2 = c^2$ ， $\triangle ABC$ 是直角三角形。

\end{answer}
